% !TeX root = ../main.tex
\part{Sản phẩm}
\section{Mô tả sản phẩm và dịch vụ}

\subsection{Tổng quan về sản phẩm và dịch vụ}
    Nhóm phát triển một nền tảng trực tuyến chuyên cung cấp dịch vụ cho thuê Gear công nghệ, bao gồm bàn phím cơ, tai nghe, tay cầm chơi game và nhiều thiết bị ngoại vi khác. Với thông điệp ``Trải nghiệm Gear xịn, không cần chi phí lớn'', nền tảng giúp sinh viên và người dùng trẻ dễ dàng tiếp cận những sản phẩm chất lượng cao mà không phải chịu áp lực tài chính khi mua mới.

    Thông qua website, khách hàng có thể tìm kiếm và so sánh thiết bị, lựa chọn thời gian thuê, đặt lịch và thanh toán trực tuyến một cách nhanh chóng, thuận tiện. Không chỉ phục vụ người thuê, nền tảng còn kết nối các cửa hàng Gear nhỏ lẻ với cộng đồng yêu công nghệ, tạo thêm một kênh kinh doanh hiệu quả. Qua đó, các đối tác có thể mở rộng tệp khách hàng, gia tăng tần suất khai thác sản phẩm và tối ưu doanh thu từ nguồn thiết bị sẵn có.

\subsection{Vấn đề của khách hàng mà sản phẩm giải quyết}
    Nhiều người dùng mong muốn trải nghiệm cảm giác gõ khác biệt của một chiếc bàn phím cơ cao cấp, chất âm sống động từ tai nghe chất lượng hay độ phản hồi chính xác của tay cầm chuyên nghiệp. Tuy nhiên, mức giá cao cùng nỗi lo mua nhầm sản phẩm không phù hợp khiến họ ngần ngại đầu tư.

    Nền tảng của nhóm giải quyết những rào cản này bằng mô hình ``thuê để trải nghiệm''. Người dùng có thể sử dụng Gear cao cấp với chi phí hợp lý, thử nghiệm sản phẩm trong điều kiện thực tế và đưa ra quyết định mua sắm tự tin hơn. Giải pháp đặc biệt phù hợp với sinh viên, game thủ trẻ, người sáng tạo nội dung và những người yêu công nghệ có nhu cầu sử dụng thiết bị trong thời gian ngắn.

    Ở phía đối tác, nhiều cửa hàng Gear nhỏ lẻ vẫn phụ thuộc vào lượng khách trực tiếp, phạm vi tiếp cận hạn chế và chưa khai thác tối đa giá trị của hàng trưng bày hoặc thiết bị sẵn có. Thông qua nền tảng, các cửa hàng có thể tiếp cận nhiều khách hàng tiềm năng hơn, tăng vòng quay khai thác sản phẩm và xây dựng thêm nguồn doanh thu từ hoạt động cho thuê.

    Nhờ đó, nền tảng tạo ra giá trị kép cho toàn bộ hệ sinh thái: khách hàng được trải nghiệm Gear chất lượng với chi phí tối ưu, trong khi các đối tác có thêm cơ hội tăng trưởng và mở rộng thị trường. Đây không chỉ là một dịch vụ cho thuê thiết bị, mà còn là cầu nối giúp trải nghiệm công nghệ cao cấp trở nên gần gũi, linh hoạt và dễ tiếp cận hơn.